\section{Contexto del problema}

Poder identificar una flor a partir de una fotografía tiene aplicaciones reales.
Una de ellas es el apoyo a plataformas de ciencia ciudadana y de monitoreo de
biodiversidad. Por ejemplo, iNaturalist ha acumulado más de 300 millones de
observaciones hechas por personas comunes y utiliza un modelo de visión por
computadora, entrenado con esas mismas observaciones, para sugerir la especie a
partir de una fotografía \autocite{inaturalist}. Sin este apoyo automático, cada
observación dependería de que una persona la revisara detenidamente, lo que
haría el proceso lento y difícil de escalar. Automatizar parte de la tarea
mejora la eficiencia y aporta en gran medida a la obtención de datos sobre
biodiversidad. Otra aplicación es la educativa, por ejemplo, en aplicaciones
para estudiantes o visitantes de jardines botánicos que permitan reconocer
flores y consultar información sobre ellas. También existe un interés
comercial: en México la floricultura es una actividad económica importante, y
el Estado de México es el principal productor de clavel y de crisantemo, con el
81.4\% del valor nacional de las flores de la temporada de otoño de 2021
\autocite{sader2022flores}. Este proyecto es una versión a menor escala de este
tipo de sistemas.

El reto es que la diversidad del mundo natural es enorme: se estima que existen
más de 5,000 especies de mamíferos, 10,000 de aves y 30,000 de peces, pero hay
más de 400,000 tipos de flores \autocite{kagglePetalsMetal}. Distinguir entre
ellas llega a requerir conocimiento especializado y, incluso para una persona
con experiencia, muchas flores son visualmente muy parecidas entre sí.
Identificarlas de forma manual es lento y no es escalable cuando se tienen miles
de fotografías, por lo que resulta valioso contar con modelos de clasificación
de imágenes que reconozcan el tipo de flor de forma automática
\autocite{ibmWhatImage}.

Este proyecto parte de la competencia \textit{Petals to the Metal}, de Kaggle,
cuyo reto es construir un modelo que identifique el tipo de flor en un conjunto
de imágenes. La competencia se limita a 104 clases de flores y consiste en un
problema de clasificación multiclase de imágenes en el que el modelo debe
identificar rasgos finos como la forma, la estructura y el color para distinguir
entre categorías cercanas.

Este problema tiene varias dificultades que lo hacen relevante desde un enfoque
de redes neuronales. Primero, las clases son heterogéneas, ya que algunas solo
tienen un subtipo de flor y otras agrupan muchas subvariedades. Además, las
imágenes no fueron capturadas en el mismo entorno, lo que les da más variedad.
Por último, el número de imágenes por clase tampoco es uniforme, por lo que,
considerando la métrica de evaluación, es importante que el modelo funcione bien
tanto en clases con pocos ejemplos como en aquellas con muchos.
